
\section{Introduction}

The computation of weight multiplicities in finite-dimensional simple representations of Lie algebras and algebraic groups is a classical problem in representation theory and algebraic combinatorics, going back to the Weyl character formula. Lusztig~\cite{lusztig1983singularities} introduced a $q$-analogue of these multiplicities, denoted by $K_{\lambda,\mu}(q)$. From a representation-theoretic point of view, the coefficients of $K_{\lambda,\mu}(q)$ record the dimensions of the graded pieces of the Brylinski--Kostant filtration on the $\mu$-weight space of the irreducible representation $L(\lambda)$~\cite{brylinski1989limits}.

In type~$A$, Lusztig's $q$-weight multiplicities are the classical Kostka--Foulkes polynomials. Lascoux and Sch\"utzenberger~\cite{lascoux1978conjecture} gave a celebrated combinatorial formula for these polynomials in terms of the \emph{charge} statistic on semistandard Young tableaux. More recently, Choi, Kim, and Lee~\cite{choilusztig} announced a combinatorial formula for $K_{\lambda,\mu}(q)$ in type~$C$, expressed through energy functions on Kirillov--Reshetikhin crystals, thereby extending the role played by charge in type~$A$.

Several other approaches to $q$-weight multiplicities are available, including Kostant's multiplicity formula~\cite{kostant1958formula}, Kazhdan--Lusztig polynomials~\cite{Lusztig1979}, Littelmann's path model~\cite{littelmann1995paths}, and Sahi's recursive algorithm~\cite{sahi2000new}. 
One of the purposes of this paper is to add a new combinatorial algorithm to this distinguished list.

%%%%%%%%%%%%%%%%%%%%%%%%%%%%%%%%%%%%%%%%%%%%%%%%%%%%%%%%%%%%%%%%%%%%%%%%
\subsection{Atomic decomposition}

Let $\Phi$ be an irreducible root system. 
We fix a system of positive roots $\Phi^+$ and denote by $\Delta$ the corresponding set of simple roots. 
Let $X$ be the corresponding weight lattice and let $X^+$ be the set of dominant weights.
Let $W_f$ and $W_a$ be the finite and affine Weyl groups attached to $\Phi$, respectively.

To each dominant weight $\lambda\in X^+$ we associate an element $\theta(\lambda)\in W_a$. Under the standard identification of elements of $W_a$ with alcoves, $\theta(\lambda)$ corresponds to the alcove obtained by translating by $\lambda$ the alcove associated with the longest element $w_0\in W_f$; see \Cref{def: theta lambda} for the precise definition.

Let $\mathcal{H}=\mathcal{H}(W_a)$ be the Hecke algebra of $W_a$. It is a $\mathbb{Z}[q^{\frac{1}{2}},q^{-\frac{1}{2}}]$-algebra with standard basis $\{\mathbf{H}_w\mid w\in W_a\}$ and canonical basis $\{\underline{\mathbf{H}}_w\mid w\in W_a\}$. For $\lambda\in X^+$ we write $\underline{\mathbf{H}}_\lambda=\underline{\mathbf{H}}_{\theta(\lambda)}$.

The spherical Hecke algebra, denoted by $\mathcal{H}^{\textbf{sph}}$, is a submodule of $\mathcal{H}$. As a $\mathbb{Z}[q^{\frac{1}{2}},q^{-\frac{1}{2}}]$-module, it is free with canonical basis $\{\underline{\mathbf{H}}_\lambda\mid \lambda\in X^+\}$. It also has a standard basis $\{\mathbf{H}_\lambda\mid \lambda\in X^+\}$. Here $\mathbf{H}_\lambda$ is not the standard Hecke basis element $\mathbf{H}_{\theta(\lambda)}$, but rather a sum of standard basis elements indexed by the double coset associated with $\theta(\lambda)$; see \cref{eq: defi standard basis} for details.

The transition coefficients between these two bases are Lusztig's $q$-weight multiplicities:
\begin{equation} \label{eq: intro A}
    \underline{\mathbf{H}}_{\lambda} =\sum_{\mu \in X^+}  K_{\lambda,\mu}(q) \mathbf{H}_\mu . 
\end{equation}

There is a third basis $\{\bfN_\lambda\mid \lambda\in X^+\}$ of $\mathcal{H}^{\textbf{sph}}$, which we call the \emph{atomic basis}. It is defined by
\begin{equation}  \label{eq: intro B}
\mathbf{N}_{\lambda}
= \sum_{\substack{\mu \in X^{+}\\ \mu \le \lambda}}
  q^{\operatorname{ht}(\lambda-\mu)}  \mathbf{H}_{\mu},
\end{equation}
where $\leq$ denotes the dominance order on $X$, and $\operatorname{ht}$ denotes the height of an element of $\mathbb{Z}\Phi$.

We may therefore write
\begin{equation} \label{eq: intro C}
    \underline{\mathbf{H}}_\lambda \,=\,\sum_{\mu\in X^{+}} a_{\mu , \lambda}(q^{\frac{1}{2}})  \bfN_\mu ,
\end{equation}
with $a_{\mu,\lambda}(q^{\frac{1}{2}})\in \mathbb{Z}[q^{\frac{1}{2}}]$. We call these polynomials \emph{atomic polynomials}. We say that $\underline{\mathbf{H}}_\lambda$ admits an atomic decomposition if
\[
    a_{\mu,\lambda}(q^{\frac{1}{2}})\in \mathbb{Z}_{\geq 0}[q^{\frac{1}{2}}]
\]
for all $\mu\in X^+$.

In type~$A$, the existence of an atomic decomposition follows from work of Lascoux~\cite{lascoux1989cyclic}, Shimozono~\cite{shimozono2001multi}, and Lecouvey--Lenart~\cite{lecouvey2021atomic}. More precisely, these sources prove the Satake counterpart of the atomic decomposition used here. The Satake transform is an isomorphism between the spherical Hecke algebra and the ring of symmetric functions; see, for instance, \cite{stembridgekostka}. Under this isomorphism, the canonical basis corresponds to Weyl characters, hence to Schur functions in type~$A$; the standard basis corresponds to Hall--Littlewood polynomials; and the atomic basis corresponds to the atoms introduced by Lascoux in~\cite{lascoux1989cyclic}.

Although these results imply, in principle, that the atomic polynomials can be computed, explicit computation remains difficult in practice. Our first main contribution is the following.

\begin{introtheorem}\label{Teo A}
    In type $A$, we construct a combinatorial algorithm for computing the atomic polynomials in \eqref{eq: intro C}. In particular, we give a new proof of the existence of an atomic decomposition in type $A$.
\end{introtheorem}

Combining \eqref{eq: intro A}, \eqref{eq: intro B}, and \eqref{eq: intro C}, one sees that the atomic polynomials determine the Kostka--Foulkes polynomials. Thus \Cref{Teo A} implies the following consequence.

\begin{introtheorem} \label{Teo B}
     In type $A$, we give a new combinatorial algorithm for computing Kostka--Foulkes polynomials.
\end{introtheorem}

\subsection{Pre-canonical bases}

The main tools in the proof of \Cref{Teo A} are the pre-canonical bases for spherical Hecke algebras introduced by Libedinsky, Patimo, and the first author~\cite{LPP}.

For an integer $i\geq 2$, let $\Phi^{\geq i}$ be the set of positive roots of height at least $i$; that is,
\[
    \Phi^{\geq i}=\{\alpha\in\Phi^+\mid \operatorname{ht}(\alpha)\geq i\}.
\]
For a subset $I\subset \Phi^+$, set $\Sigma_I=\sum_{\alpha\in I}\alpha$.

A weight $\lambda$ is called singular if there exists a reflection $t\in W_f$ such that $t(\lambda+\rho)=\lambda+\rho$, where $\rho$ denotes the half-sum of the positive roots. Otherwise, $\lambda$ is called regular. If $\lambda\in X$ is regular, let $w_\lambda\in W_f$ be the unique element such that $w_\lambda\cdot \lambda\in X^+$, where the dot action is defined by
\[
    w\cdot \lambda=w(\lambda+\rho)-\rho.
\]
For $\lambda\in X$, define
\begin{equation}  \label{tilde de schur}
    \widetilde{\underline{\mathbf{H}}}_\lambda = \begin{cases}
        (-1)^{\ell (w_{\lambda})} \underline{\mathbf{H}}_{w_{\lambda}\cdot \lambda}, & \mbox{if } \lambda \mbox{ is regular;}\\
        0, & \mbox{if } \lambda \mbox{ is singular.}
    \end{cases}
\end{equation}

\begin{definition}\rm \label{defi pre-canonical}
For $i\geq 2$, the $i$-th pre-canonical basis $\mathcal{N}^{i}=\{\bfN_{\lambda }^i\mid \lambda\in X^{+}\}$ is defined by
\begin{equation}  \label{def: pre-canonical}
    \bfN_{\lambda}^{i} = \sum_{I\subset \Phi^{\geq i}} (-q)^{|I|} \widetilde{\underline{\mathbf{H}}}_{\lambda - \Sigma_I}.
\end{equation}  
\end{definition}

The motivation for this definition is the anti-atomic formula~\cite[Theorem 3.1]{LPP}, which asserts that $\bfN_\lambda=\bfN_\lambda^2$ for all $\lambda\in X^+$. Thus the second pre-canonical basis is precisely the atomic basis. This formula expresses atomic basis elements in terms of the canonical basis, and it holds in every affine type.

At the other end of the interpolation, if $h$ is the height of the highest root in $\Phi$, then $\bfN_\lambda^{m}=\underline{\mathbf{H}}_\lambda$ for all $m>h$ and all $\lambda\in X^+$. Hence the pre-canonical bases interpolate between the atomic basis and the canonical basis. %If one also takes \eqref{eq: intro B} into account, they may be viewed as interpolating from the standard basis to the canonical basis through the atomic basis.
Together with \eqref{eq: intro B}, this shows that the sequence of pre-canonical bases interpolates from the standard basis to the canonical basis, passing through the atomic basis at $i=2$.

Based on extensive computations, the authors of~\cite{LPP} proposed the following positivity conjecture.

\begin{introconjecture}   \label{Conj A}
    In type $A$, each element of $\mathcal{N}^{i+1}$ expands positively in the basis $\mathcal{N}^i$.
    
    Equivalently, if
    \begin{equation}
        \bfN_{\lambda}^{i+1} =\sum_{\mu\in X^+} a_{\mu , \lambda}^i (q)  \bfN_{\mu}^i,
    \end{equation}
    then $a_{\mu , \lambda}^i(q)\in\mathbb{Z}_{\geq 0}[q]$ for all $\lambda,\mu\in X^+$ and all $i\geq 2$.
\end{introconjecture}

Our second main contribution is the following theorem.

\begin{introtheorem} \label{Teo C}
    In type $A$, we give a combinatorial algorithm for computing the polynomials $a_{\mu , \lambda}^i(q)$. Using this algorithm, we prove that $a_{\mu , \lambda}^i(q)\in\mathbb{Z}_{\geq 0}[q]$. In particular, we prove \Cref{Conj A}.
\end{introtheorem}

We remark that \Cref{Teo A} and \Cref{Teo B} follow readily from \Cref{Teo C}.

We now outline the proof of \Cref{Teo C}. Fix a root system $\Phi$ of type $A_n$, and let $\Delta=\{\alpha_1,\ldots,\alpha_n\}$ be the set of simple roots. For $1\leq i\leq j\leq n$, set
\begin{equation}
     \alpha_{i,j} =\sum_{k=i}^j \alpha_k \in \Phi^+.
\end{equation}
Thus $\operatorname{ht}(\alpha_{i,j})=j-i+1$. We order $\Phi^{\geq 2}$ by declaring that $\alpha_{i,j}\preceq \alpha_{a,b}$ if either
\begin{itemize}
    \item $\operatorname{ht}(\alpha_{i,j})<\operatorname{ht}(\alpha_{a,b})$, or
    \item $\operatorname{ht}(\alpha_{i,j})=\operatorname{ht}(\alpha_{a,b})$ and $i\leq a$.
\end{itemize}

Define
\[
    \Phi^{\succeq \alpha_{i,j}}=\{\alpha\in\Phi^{\geq 2}\mid \alpha\succeq \alpha_{i,j}\},
    \qquad
    \Phi^{\succ \alpha_{i,j}}=\{\alpha\in\Phi^{\geq 2}\mid \alpha\succ \alpha_{i,j}\}.
\]
We also set
\begin{equation}
    \bfM_{\lambda}^{\succeq \alpha_{i,j}}  =  \sum_{I\subset \Phi^{\succeq \alpha_{i,j} }} (-q)^{|I|} \widetilde{\underline{\mathbf{H}}}_{\lambda - \Sigma_I} 
    \qquad \mbox{and} \qquad 
    \bfM_{\lambda}^{\succ \alpha_{i,j}}  =  \sum_{I\subset \Phi^{\succ \alpha_{i,j} }} (-q)^{|I|} \widetilde{\underline{\mathbf{H}}}_{\lambda - \Sigma_I}.
\end{equation}

Ideally, one would like to express $\bfM_{\lambda}^{\succ \alpha_{i,j}}$ as a positive linear combination of the elements
\[
    \{\bfM_{\mu}^{\succeq \alpha_{i,j}}\mid \mu\in X^+\}.
\]
Iterating such an expression would prove \Cref{Teo C}, since
\begin{equation}
    \Phi^{\succ \alpha_{n-i+1,n}} =\Phi^{\geq i+1}
    \qquad \mbox{and} \qquad
    \Phi^{\succeq \alpha_{1,i}} = \Phi^{\geq i};
\end{equation}
hence
\begin{equation}\label{eq: M es N intro}
    \bfN_{\lambda}^{i+1} = \bfM_{\lambda}^{\succ \alpha_{n-i+1,n}} 
    \qquad \mbox{and} \qquad 
    \bfN_{\lambda}^i = \bfM_{\lambda}^{\succeq \alpha_{1,i}}, \mbox{ for all } 2\leq i\leq n.
\end{equation}

However, this direct strategy does not work: such a decomposition is difficult to obtain and, more importantly, it is not positive in general. The key point of the paper is to replace this unattainable decomposition by a more flexible positive statement.

\begin{introtheorem} \label{Teo D}
Let $\alpha_{i,j} \in \Phi^{\geq 2}$ and $\lambda \in X^+$. Then there exist polynomials $b_{\mu, \lambda }^k(q) \in \mathbb{Z}_{\geq 0}[q]$ such that
\begin{equation} \label{eq: intro descomposition}
    \bfM_{\lambda}^{\succ \alpha_{i,j} } = \sum_{k=0}^{i-1}  \sum_{\substack{\mu \in X^{+}\\ \mu \le \lambda}} b_{\mu, \lambda }^k(q) \bfM_{ \mu }^{\succeq \alpha_{i-k,j-k}}.
\end{equation}
Moreover, we give a combinatorial algorithm for computing the polynomials $b_{\mu,\lambda}^k(q)$.
\end{introtheorem}

Repeatedly applying \Cref{Teo D} yields the following strengthening of \Cref{Teo C}.

\begin{introtheorem}  \label{Teo E}
    Let $\alpha_{i,j}\in\Phi^{\geq 2}$, let $h=j-i+1$, and let $\lambda\in X^+$. Then $\bfM_{\lambda}^{\succ \alpha_{i,j}}$ can be written as a positive linear combination of elements of the $h$-th pre-canonical basis.
    
    In particular, $\bfN^{h+1}_\lambda=\bfM_{\lambda}^{\succ \alpha_{n-h+1,n}}$ can be written as a positive linear combination of the $h$-th pre-canonical basis, which is precisely the assertion of \Cref{Teo C}.
\end{introtheorem}

%%%%%%%%%%%%%%%%%%%%%%%%%%%%%%%%%%%%%%%%%%%%%%%%%%%%%%%%%%%%%%%%%%%%%%%%
\subsection{An illustrative example}

We illustrate \Cref{Teo D} with a concrete example. The purpose of the example is not to present the full technical construction, but rather to show the kind of combinatorics that underlies the proof. The definitions of the operators used below are given later in \Cref{def. p and q,def: words over weights,def: v negrita,def: P negrita}.

Let $n=14$, let $\alpha_{i,j}=\alpha_{8,10}$, and let $h=j-i+1=3$. Consider the dominant weight
\begin{equation}
\lambda = [0,1,0,2,1,1,0,0,0,0,0,0,1,1]_{\varpi},
\end{equation}
where $\lambda=[\lambda_1,\ldots,\lambda_{14}]_{\varpi}$ means that $\lambda=\sum_{k=1}^{14}\lambda_k\varpi_k$, and $\{\varpi_k\}_{1\leq k\leq 14}$ are the fundamental weights.

Our goal is to express $\bfM_{\lambda}^{\succ \alpha_{8,10}}$ as a positive linear combination of terms in
\begin{equation}
 \mathcal{M}\coloneqq \{\bfM_{\mu}^{\succeq \alpha_{8-k,10-k}}\mid \mu\in X^+,\, \mu\leq\lambda, \mbox{ and } 1\leq k\leq 7\}.
\end{equation}

The starting point of our analysis is \Cref{prop: second version refinado}. 
We refer to the formula in that proposition as the \emph{Second Inverse Decomposition} (SID). 
This terminology is justified by the fact that the SID is a refinement of the so-called \emph{First Inverse Decomposition} (FID) in \Cref{prop: second version}, which yields a preliminary decomposition in which some of the weights involved are not necessarily dominant.


In the present example, the SID gives
\begin{equation} \label{eq:ejemplo-intro}
\begin{array}{rlll}
        & \quad\bfM_\lambda^{\succeq \alpha_{8,10}}  + & &  \\
        &   &  &  \\
   &  q^9 \bfM_{\mathbf{P}_{10,3}(\lambda)}^{\succ \alpha_{8,10}} + & & \mathbf{P}\mbox{-term} \\
   \bfM_\lambda^{\succ \alpha_{8,10}} =  & &  &    \\
     &q^{4}\!\left(\bfM_{v_{10,3}^1(\lambda)}^{\succ \alpha_{8,10}} 
- q\, \bfM_{v_{10,3}^1(\lambda)-\alpha_{7,9}}^{\succ \alpha_{8,10}}\right) + & & v^1\mbox{-term}  \\
& & &  \\
     & q^{8}\!\left(\bfM_{\mathbf{v}_{10,3}(\lambda)}^{\succ \alpha_{8,10}} 
- q^5\, \bfM_{\mathbf{P}_{8,3}(\mathbf{v}_{10,3}(\lambda))}^{\succ \alpha_{8,10}}\right).  & & \mathbf{v}\mbox{-term and } \mathbf{Pv}\mbox{-term}  \\
\end{array}
\end{equation}

The weights appearing in this formula are
\begin{equation}
\begin{array}{r}
\mathbf{P}_{10,3}(\lambda)= [1,1,1,1,0,0,0,0,0,0,1,0,1,1]_{\varpi},\\[3pt]
 v_{10,3}^1(\lambda)= [0,1,0,3,0,1,0,0,0,0,0,0,1,0]_{\varpi}, \\[3pt]
\mathbf{v}_{10,3}(\lambda)= [0,1,1,2,0,1,0,0,0,0,0,0,0,1]_{\varpi}, \\[3pt]
\mathbf{P}_{8,3}(\mathbf{v}_{10,3}(\lambda)) = [1, 1, 1, 1, 0, 0, 0, 0, 1, 0, 0, 0, 0, 1]_{\varpi}.
\end{array}
\end{equation}

We now explain how the weights and exponents in \eqref{eq:ejemplo-intro} arise. We identify dominant weights with integer partitions, and integer partitions with Young diagrams. We use partitions with at most $15$ parts. Thus, if $\mu=(\mu_1,\ldots,\mu_{15})$, then the corresponding dominant weight is
\begin{equation}
[\mu_1-\mu_2,\mu_2-\mu_3,\ldots,\mu_{14}-\mu_{15}]_{\varpi}.
\end{equation}

\begin{itemize}
    \item The term $\bfM_\lambda^{\succeq \alpha_{8,10}}$ always appears with coefficient $1$.

\item Consider the $\mathbf{P}$-term $q^9\bfM_{\mathbf{P}_{10,3}(\lambda)}^{\succ \alpha_{8,10}}$.
 
In the partition model, subtracting a root $\alpha_{i,j}$ amounts to moving one box from row $i$ to row $j+1$. The resulting diagram need not be a partition: the box may fail to be removable from row $i$, or its new position in row $j+1$ may fail to be addable. 

The construction of $\mathbf{P}_{10,3}(\lambda)$ is recursive. It is obtained as a composition of certain operators, which we denote by $P$ (in non-boldface), to distinguish them from the boldface objects $\mathbf{P}_{j,h}(\lambda)$.

We begin by considering the  root $\alpha_{8,10}$ of height $h=3$. If row $8$ has a removable box, we set
\begin{equation}
P_{10,3}(\lambda)=\lambda-\alpha_{8,10}.
\end{equation}
If not, we check row $8-h=5$; if row $5$ has a removable box, we set
\begin{equation}
P_{10,3}(\lambda)=\lambda-\alpha_{5,7}-\alpha_{8,10}.
\end{equation}
If this also fails, we continue with row $5-h=2$, and so on. In general, we look for a row $\varrho\equiv i\pmod h$ from which a box can be moved, and define
\begin{equation}
P_{j,h}(\lambda)=\lambda-\alpha_{\varrho,j}.
\end{equation}
If no such row exists, both $P_{j,h}(\lambda)$ and $\mathbf{P}_{j,h}(\lambda)$ are undefined.

In our example,
\begin{equation}
P_{10,3}(\lambda)=\lambda-\alpha_{5,7}-\alpha_{8,10}.
\end{equation}
We then check whether this weight is dominant. If it is, we set $\mathbf{P}_{10,3}(\lambda)=P_{10,3}(\lambda)$. If it is not, we repeat the same procedure starting from $\alpha_{7,9}=\alpha_{8-1,10-1}$; this gives the operator $P_{9,3}$. If $P_{9,3}P_{10,3}(\lambda)$ is defined but not dominant, we continue with the root $\alpha_{6,8}=\alpha_{8-2,10-2}$, and so on.

The process stops either when the resulting weight becomes dominant, in which case this weight is $\mathbf{P}_{10,3}(\lambda)$, or when the process terminates unsuccessfully, in which case $\mathbf{P}_{10,3}(\lambda)$ is undefined. This is illustrated in \Cref{fig:introA}.

Finally, passing from $\lambda$ to $\mathbf{P}_{10,3}(\lambda)$ requires subtracting nine roots of height $3$. This explains the coefficient $q^9$ in front of $\bfM_{\mathbf{P}_{10,3}(\lambda)}^{\succ \alpha_{8,10}}$ in \eqref{eq:ejemplo-intro}.
\begin{figure}[h]
    \centering


\Yboxdim{1cm}
\begin{tikzpicture}[scale=.21]
\tikzset{>={Stealth[length=1.4mm, width=1.1mm]}}

\Yfillopacity{0.3}
\Ylinethick{0.3pt}
\Ylinecolour{blue}
\Yfillcolour{cyan}

% --- Diagram 1 ---
\begin{scope}[local bounding box=Y0]
  \tyng(0cm,3cm,7^2,6^2,4,3,2^7,1)
  \draw[->,red,line width=0.4mm, line cap=round](2.5,-3.5)--(3.5,-3.5)%
    to node[auto,swap]{} (3.5,-6.5)--(2.3,-6.5);
  \draw[->,red,line width=0.4mm, line cap=round](4.5,-0.5)--(5.5,-0.5)%
    to node[auto,swap]{} (5.5,-3.5)--(4,-3.5);
\end{scope}

% --- Diagram 2 ---
\begin{scope}[xshift=9cm,local bounding box=Y1]
  \tyng(0cm,3cm,7^2,6^2,3^2,2^4,3,2^2,1)
  \draw[->,red,line width=0.4mm, line cap=round](2.5,-2.5)--(3.5,-2.5)%
    to node[auto,swap]{} (3.5,-5.5)--(2.3,-5.5);
  \draw[->,red,line width=0.4mm, line cap=round](6.5,0.5)--(7,0.5)%
    to node[auto,swap]{} (7,-2.5)--(4,-2.5);
\end{scope}

% --- Diagram 3 ---
\begin{scope}[xshift=18cm,local bounding box=Y2]
  \tyng(0cm,3cm,7^2,6,5,3^2,2^3,3^2,2^2,1)
  \draw[->,red,line width=0.4mm, line cap=round](3.5,-1.5)--(4,-1.5)%
    to node[auto,swap]{} (4,-4.5)--(2.8,-4.5);
\end{scope}

% --- Diagram 4 ---
\begin{scope}[xshift=27cm,local bounding box=Y3]
  \tyng(0cm,3cm,7^2,6,5,3,2^3,3^3,2^2,1)
  \draw[->,red,line width=0.4mm, line cap=round](3.5,-0.5)--(4,-0.5)%
    to node[auto,swap]{} (4,-3.5)--(2.5,-3.5);
\end{scope}


% --- Diagram 5 ---
\begin{scope}[xshift=36cm,local bounding box=Y4]
  \tyng(0cm,3cm,7^2,6,5,2^3,3^4,2^2,1)
  \draw[->,red,line width=0.4mm, line cap=round](5.5,0.5)--(6,0.5)%
    to node[auto,swap]{} (6,-2.5)--(2.8,-2.5);
\end{scope}


% --- Diagram 6 ---
\begin{scope}[xshift=45cm,local bounding box=Y5]
  \tyng(0cm,3cm,7^2,6,4,2^2,3^5,2^2,1)
  \draw[->,red,line width=0.4mm, line cap=round](6.5,1.5)--(7,1.5)%
    to node[auto,swap]{} (7,-1.5)--(2.8,-1.5);
\end{scope}


% --- Diagram 7 ---
\begin{scope}[xshift=54cm,local bounding box=Y6]
  \tyng(0cm,3cm,7^2,5,4,2,3^6,2^2,1)
  \draw[->,red,line width=0.4mm, line cap=round](7.5,2.5)--(8,2.5)%
    to node[auto,swap]{} (8,-0.5)--(2.8,-0.5);
\end{scope}


% --- Diagram 8 ---
\begin{scope}[xshift=63cm,local bounding box=Y7]
  \tyng(0cm,3cm,7,6,5,4,3^7,2^2,1)
\end{scope}



% Símbolos centrados entre cada par de tableaux
\foreach \A/\B/\Sym in {
  Y0/Y1/{$P_{10,3} $},
  Y1/Y2/{$P_{9,3} $},
  Y2/Y3/{$P_{8,3} $},
  Y3/Y4/{$P_{7,3} $},
  Y4/Y5/{$P_{6,3} $},
  Y5/Y6/{$P_{5,3} $},
  Y6/Y7/{$P_{4,3} $}
}{
  \node[font=\small,anchor=north,yshift=-3pt]
       at ($(\A.south)!0.5!(\B.south)$) {\Sym};
}


\end{tikzpicture}
\caption{The leftmost diagram represents the partition $\lambda$, while the rightmost one corresponds to $\mathbf{P}_{10,3}(\lambda)$. 
Each intermediate diagram is obtained from the previous one by applying the operator displayed between them.}
    \label{fig:introA}
\end{figure}

\item We now turn to the $v^1$-term $v^1_{10,3}(\lambda)$.

To obtain this weight, we first repeat the initial step of the previous construction, namely we form $P_{10,3}(\lambda)$. We then move the box in row $11$ by subtracting $\alpha_{11,12}$, a root of height $h-1=2$. If the moved box reaches an addable position, the process stops. Otherwise, we subtract $\alpha_{13,14}=\alpha_{11+2,12+2}$ and continue until either the box reaches an addable position or the resulting weight ceases to exist. This operator is denoted by $Q_{10,3}$, and we define
\begin{equation}
v^1_{10,3}(\lambda)=Q_{10,3}(P_{10,3}(\lambda)).
\end{equation}
The construction is illustrated in \Cref{fig:introB}. The superscript $1$ in $v^1_{10,3}$ indicates that we have used one $P$-operator and one $Q$-operator. The exponent $4$ records the four roots that must be subtracted from $\lambda$ to obtain $v^1_{10,3}(\lambda)$.
\begin{figure}[h]
    \centering
    \Yboxdim{1cm}
\begin{tikzpicture}[scale=.21]
\tikzset{>={Stealth[length=1.4mm, width=1.1mm]}}

\Yfillopacity{0.3}
\Ylinethick{0.3pt}
\Ylinecolour{blue}
\Yfillcolour{cyan}

% --- Diagram 1 ---
\begin{scope}[local bounding box=Y0]
  \tyng(0cm,3cm,7^2,6^2,4,3,2^7,1)
  \tikzset{>={Stealth[length=1.4mm, width=1.1mm]}}
  \draw[->,red,line width=0.4mm, line cap=round](2.5,-3.5)--(3.5,-3.5)%
    to node[auto,swap]{}%
    (3.5,-6.5)--(2.3,-6.5);
  \draw[->,red,line width=0.4mm, line cap=round](4.5,-0.5)--(5.5,-0.5)%
    to node[auto,swap]{}%
    (5.5,-3.5)--(4,-3.5);
\end{scope}



% --- Diagram 2 ---
\begin{scope}[xshift=10.5cm, local bounding box=Y1]
  \tyng(0cm,3cm,7^2,6^2,3^2,2^4,3,2^2,1)
  \draw[->,blue,line width=0.4mm, line cap=round](3.5,-6.5)--(4,-6.5)%
    to node[auto,swap]{}%
    (4,-8.5)--(2.5,-8.5);
  \draw[->,blue,line width=0.4mm, line cap=round](2.5,-8.5)--(2.5,-8.5)%
    to node[auto,swap]{}%
    (2.5,-10.5)--(0.5,-10.5);
\end{scope}


% --- Diagram 3 ---
\begin{scope}[xshift=21cm, local bounding box=Y2]
  \tyng(0cm,3cm,7^2,6^2,3^2,2^7,1^2)
\end{scope}



% Símbolos centrados entre cada par de tableaux
\foreach \A/\B/\Sym in {
  Y0/Y1/{$P_{10,3} $},
  Y1/Y2/{$Q_{10,3} $}
}{
  \node[font=\small,anchor=north,yshift=-3pt]
       at ($(\A.south)!0.5!(\B.south)$) {\Sym};
}

\end{tikzpicture}

    \caption{The leftmost diagram represents the partition $\lambda$, while the rightmost one corresponds to $v_{10,3}^1(\lambda)$.}
    \label{fig:introB}
\end{figure}

The weight $v_{10,3}^1(\lambda)-\alpha_{7,9}$ also appears in the third row of the right-hand side of \eqref{eq:ejemplo-intro}. This weight need not be dominant. It always occurs with an additional factor of $q$. Together with the external factor $q^4$, this reflects the fact that five roots must be subtracted from $\lambda$ to obtain $v_{10,3}^1(\lambda)-\alpha_{7,9}$.

\item We now consider the $\mathbf{v}$-term $\mathbf{v}_{10,3}(\lambda)$.

In this case, we repeat the first two steps in the construction of $\mathbf{P}_{10,3}(\lambda)$; in general, one repeats $h-1$ steps. Thus we begin with $P_{9,3}P_{10,3}(\lambda)$. We then apply $Q$-operators until the weight becomes dominant: in this example, we apply $Q_{9,3}$, then $Q_{10,3}$, then $Q_{11,3}$, and so on. The process is shown in \Cref{fig:introC}. As before, the exponent $q^8$ records the eight roots subtracted during the process.

\item Finally, consider the $\mathbf{Pv}$-term $\mathbf{P}_{8,3}(\mathbf{v}_{10,3}(\lambda))$.

Starting from $\mathbf{v}_{10,3}(\lambda)$, we proceed as in the construction of the $\mathbf{P}$-term, but now with the root $\alpha_{8,10}=\alpha_{i,j}$ replaced by $\alpha_{6,8}=\alpha_{i-h+1,j-h+1}$. This explains the occurrence of $\mathbf{P}_{8,3}(\mathbf{v}_{10,3}(\lambda))$ and the coefficient $q^{13}$.
\end{itemize}
\begin{figure}
    \centering
    \Yboxdim{1cm}
\begin{tikzpicture}[scale=.21]
\tikzset{>={Stealth[length=1.4mm, width=1.1mm]}}

\Yfillopacity{0.3}
\Ylinethick{0.3pt}
\Ylinecolour{blue}
\Yfillcolour{cyan}


% --- Diagram 1 ---
\begin{scope}[local bounding box=Y0]
  \tyng(0cm,3cm,7^2,6^2,4,3,2^7,1)
  \tikzset{>={Stealth[length=1.4mm, width=1.1mm]}}
  \draw[->,red,line width=0.4mm, line cap=round](2.5,-3.5)--(3.5,-3.5)%
    to node[auto,swap]{}%
    (3.5,-6.5)--(2.3,-6.5);
  \draw[->,red,line width=0.4mm, line cap=round](4.5,-0.5)--(5.5,-0.5)%
    to node[auto,swap]{}%
    (5.5,-3.5)--(4,-3.5);
\end{scope}



% --- Diagram 2 ---
\begin{scope}[xshift=10.5cm, local bounding box=Y1]
  \tyng(0cm,3cm,7^2,6^2,3^2,2^4,3,2^2,1)
  \draw[->,red,line width=0.4mm, line cap=round](2.5,-2.5)--(3.5,-2.5)%
    to node[auto,swap]{}%
    (3.5,-5.5)--(2.3,-5.5);
  \draw[->,red,line width=0.4mm, line cap=round](6.5,0.5)--(7,0.5)%
    to node[auto,swap]{}%
    (7,-2.5)--(4,-2.5);
\end{scope}


% --- Diagram 3 ---
\begin{scope}[xshift=22cm, local bounding box=Y2]
  \tyng(0cm,3cm,7^2,6,5,3^2,2^3,3^2,2^2,1)
  \draw[->,blue,line width=0.4mm, line cap=round](3.5,-5.5)--(4,-5.5)%
    to node[auto,swap]{}%
    (4,-7.5)--(2.5,-7.5);
\end{scope}


% --- Diagram 4 ---
\begin{scope}[xshift=33cm, local bounding box=Y3]
  \tyng(0cm,3cm,7^2,6,5,3^2,2^4,3^2,2^1,1)
  \draw[->,blue,line width=0.4mm, line cap=round](3.5,-6.5)--(4,-6.5)%
    to node[auto,swap]{}%
    (4,-8.5)--(2.5,-8.5);
\end{scope}


% --- Diagram 5 ---
\begin{scope}[xshift=44cm, local bounding box=Y4]
  \tyng(0cm,3cm,7^2,6,5,3^2,2^5,3^2,1)
  \draw[->,blue,line width=0.4mm, line cap=round](3.5,-7.5)--(4,-7.5)%
    to node[auto,swap]{}%
    (4,-9.5)--(1.5,-9.5);
\end{scope}


% --- Diagram 6 ---
\begin{scope}[xshift=55cm, local bounding box=Y5]
  \tyng(0cm,3cm,7^2,6,5,3^2,2^6,3,2)
  \draw[->,blue,line width=0.4mm, line cap=round](3.5,-8.5)--(4,-8.5)%
    to node[auto,swap]{}%
    (4,-10.5)--(0.5,-10.5);
\end{scope}


% --- Diagram 7 ---
\begin{scope}[xshift=66cm, local bounding box=Y6]
  \tyng(0cm,3cm,7^2,6,5,3^2,2^8,1)
\end{scope}



% Símbolos centrados entre cada par de tableaux
\foreach \A/\B/\Sym in {
  Y0/Y1/{$P_{10,3} $},
  Y1/Y2/{$P_{9,3} $},
  Y2/Y3/{$Q_{9,3} $},
  Y3/Y4/{$Q_{10,3} $},
  Y4/Y5/{$Q_{11,3} $},
  Y5/Y6/{$Q_{12,3} $}
}{
  \node[font=\small,anchor=north,yshift=-3pt]
       at ($(\A.south)!0.5!(\B.south)$) {\Sym};
}


\end{tikzpicture}

    \caption{The leftmost diagram represents the partition $\lambda$, while the rightmost one corresponds to $\mathbf{v}_{10,3}(\lambda) $.}
    \label{fig:introC}
\end{figure}

The existence of the weights above depends strongly on $\lambda$. A useful feature of the SID is that, whenever one of these weights is undefined, the corresponding $\bfM$-term is simply omitted and the formula remains valid. For example, if $\lambda$ is the zero weight or a fundamental weight, then none of the auxiliary weights appearing above is defined, and the SID reduces to
\[
    \bfM_{\lambda}^{\succ\alpha_{i,j}}=\bfM_{\lambda}^{\succeq\alpha_{i,j}}.
\]

We now explain how positivity is extracted from \eqref{eq:ejemplo-intro}.

\begin{itemize}
    \item The element $\bfM_{\lambda}^{\succeq\alpha_{8,10}}$ already belongs to $\mathcal{M}$.

    \item The $\mathbf{P}$-term $\bfM_{\mathbf{P}_{10,3}(\lambda)}^{\succ \alpha_{8,10}}$ is handled by induction on the dominance order, since $\mathbf{P}_{10,3}(\lambda)<\lambda$ by construction.

    \item The difference
    \[
        \bfM_{v_{10,3}^1(\lambda)}^{\succ \alpha_{8,10}}
        -q\,\bfM_{v_{10,3}^1(\lambda)-\alpha_{7,9}}^{\succ \alpha_{8,10}}
    \]
    is referred to as a \emph{good pair} in \Cref{section: positivity}.
    %is what we call a \emph{good pair} in \Cref{section: positivity}. 
    In general, the SID may produce $h-2$ good pairs. In this example there is only one, since $h-2=1$.

    The weights appearing in good pairs have the form $v_{j,h}^k(\lambda)$ and $v_{j,h}^k(\lambda)-\alpha_{i-k,j-k}$, with $1\leq k\leq h-2$. The operators $v_{j,h}^k$ are natural generalizations of the operator $v_{10,3}^1$ used above; see \Cref{def: words over weights}. By \Cref{prop: caso bueno}, each good pair admits a positive combinatorial expansion in terms of elements of the form $\bfM^{\succeq\alpha_{i-k,j-k}}$.

    In the present example, $k=1$, and we have
    \begin{equation} \label{eq:intro-expansion-good}
    \bfM_{v_{10,3}^1(\lambda)}^{\succ \alpha_{8,10}} 
    - q\, \bfM_{\,v_{10,3}^1(\lambda)-\alpha_{7,9}}^{\succ \alpha_{8,10}}
    = 
    \sum_{m=1}^{\infty} q^{e_m}\,
    \bfM_{(v^1_{10,3})^{m} (\lambda)}^{\succeq \alpha_{7,9}},
\end{equation}
where $e_m$ is the number of positive roots that must be subtracted from $v^1_{10,3}(\lambda)$ in order to obtain $(v^1_{10,3})^m(\lambda)$ and $(v_{10,3}^{1})^m$ denotes the $m$-fold composition of $v_{10,3}^1$.

With the convention that undefined weights contribute zero, the apparently infinite sum in \eqref{eq:intro-expansion-good} is finite: each application of $v^1_{10,3}$ strictly decreases the weight in the dominance order, and there are only finitely many dominant weights below a fixed dominant weight. In this example,
\begin{equation} \label{eq:intro-expansion-good-concrete}
    \bfM_{v_{10,3}^1(\lambda)}^{\succ \alpha_{8,10}} 
    - q\, \bfM_{\,v_{10,3}^1(\lambda)-\alpha_{7,9}}^{\succ \alpha_{8,10}}
    = 
    \bfM_{v^1_{10,3}(\lambda)}^{\succeq \alpha_{7,9}},
\end{equation}
since $(v_{10,3}^1)^m(\lambda)$ is undefined for all $m\geq 2$.

\item Finally, we consider the difference
\[
    \bfM_{\mathbf{v}_{10,3}(\lambda)}^{\succ \alpha_{8,10}}
    -q^5\,\bfM_{\mathbf{P}_{8,3}(\mathbf{v}_{10,3}(\lambda))}^{\succ \alpha_{8,10}}.
\]
%This is what we call a \emph{bad pair}.
This is referred to as a \emph{bad pair}.
In general, at most one bad pair appears in the SID. Unlike good pairs, bad pairs do not admit an immediate positive expansion, which explains the terminology.

The two weights appearing in a bad pair are decomposed further by applying the SID recursively.
An induction argument, together with commutation rules for the operators on weights, shows that the bad pair can nevertheless be rewritten as a positive linear combination of elements of the form $\bfM^{\succeq\alpha_{i-k,j-k}}$ for $1\leq k\leq i-1$. 
The required commutation rules are introduced in \Cref{section COM RUl}, and the induction is carried out in \Cref{section proof of positivity}.
\end{itemize}

%%%%%%%%%%%%%%%%%%%%%%%%%%%%%%%%%%%%%%%%%%%%%%%%%%%%%%%%%%%%%%%%%%%%%%
\subsection{Symmetric functions}
%%%%%%%%%%%%%%%%%%%%%%%%%%%%%%%%%%%%%%%%%%%%%%%%%%%%%%%%%%%%%%%%%%%%%%

The results of this paper can also be viewed through the language of symmetric functions. As mentioned above, the Satake isomorphism relates the spherical Hecke algebra to the algebra of symmetric functions. We now spell out the symmetric-function analogue of the pre-canonical bases.

As before, we identify dominant weights in type $A$ with partitions. Let $s_\lambda$ denote the Schur function indexed by a partition $\lambda$. We extend Schur functions to arbitrary elements of $\mathbb{Z}^{\ell}$ by the usual straightening rule: for $\gamma\in\mathbb{Z}^{\ell}$, set
\begin{equation}\label{eq:straighten}
  s_\gamma =
  \begin{cases}
    \operatorname{sgn}(\gamma+\rho)\,
    s_{\operatorname{sort}(\gamma+\rho)-\rho},
      & \text{if $\gamma+\rho$ has distinct nonnegative parts,} \\[4pt]
    0, & \text{otherwise.}
  \end{cases}
\end{equation}
Here $\rho=(\ell-1,\ell-2,\dots,0)$, $\operatorname{sort}(\beta)$ denotes the weakly decreasing rearrangement of $\beta$, and $\operatorname{sgn}(\beta)$ is the sign of the shortest permutation sending $\beta$ to $\operatorname{sort}(\beta)$.

This straightening rule is the Schur-function counterpart of \eqref{tilde de schur}. The analogue of \Cref{defi pre-canonical} is the following.

\begin{definition}\rm
The $k$-th pre-canonical Schur function associated with $\lambda$ is
\begin{equation}\label{eq:preSchur}
  s_{\lambda}^{k}
  =  \prod_{\alpha \in \Phi^{\ge k}} (1 - q\,L_{\alpha})\, s_{\lambda},
\end{equation}
where the lowering operator $L_\alpha$ acts on subscripts by
\[
    L_{\alpha}s_{\gamma}=s_{\gamma-(\epsilon_i-\epsilon_{j+1})}
    \qquad \mbox{if } \alpha=\alpha_{i,j}.
\]
\end{definition}

In \eqref{eq:preSchur}, the product is first expanded and then applied to $s_\lambda$. Thus
\begin{equation}\label{eq:preSchurA}
  s_{\lambda}^{k}
  = \sum_{I\subset \Phi^{\geq k}} (-q)^{|I|} s_{\lambda-\Sigma_I},  
\end{equation}
which is the image of \Cref{defi pre-canonical} under the heuristic replacement of canonical basis elements by Schur functions.

The formulation in terms of lowering operators highlights a parallel with Catalan symmetric functions. These are defined by
\begin{equation}
\mathbf{H}_{A,\lambda}
  = \prod_{\alpha\in A} (1 - q R_{\alpha})^{-1} s_{\lambda},
\end{equation}
where $A$ is a root ideal and the raising operator $R_\alpha$ acts by
\[
    R_{\alpha}s_\gamma=s_{\gamma+(\epsilon_i-\epsilon_{j+1})}
    \qquad \mbox{if } \alpha=\alpha_{i,j}.
\]
For our purposes, the relevant point is that each set $\Phi^{\geq k}$ is a root ideal.

Catalan symmetric functions were introduced independently in \cite{chen2010skew} and \cite{panyushev2010generalised}. They interpolate between modified Hall--Littlewood polynomials, when $A=\Phi^+$, and Schur functions, when $A=\varnothing$.

Similarly, after applying the Satake transform and then dualizing with respect to the Hall inner product, the standard and canonical bases of $\mathcal{H}^{\mathbf{sph}}$ correspond to modified Hall--Littlewood and Schur functions, respectively. It is natural to ask whether the pre-canonical Schur functions and the Catalan symmetric functions attached to the ideals $\Phi^{\geq k}$ are dual bases; that is, whether
\begin{equation}
    \langle s_{\lambda}^k , \mathbf{H}_{\Phi^{\ge k},\mu} \rangle = \delta_{\lambda,\mu}.
\end{equation}
This is not the case. For example, when $n=4$ and $k=2$,
\begin{equation}
    s^2_{(2,1,1)} 
      = s_{(2,1,1)} + (-q^{2}-q)\, s_{(1,1,1,1)},
\end{equation}
whereas
\begin{equation}
    \mathbf{H}_{\Phi^{\ge 2}, (1,1,1,1)} 
      = s_{(1,1,1,1)} + q\,s_{(2,1,1)} + q^2 s_{(2,2)}.
\end{equation}
Hence
\begin{equation}
  \langle s^2_{(2,1,1)} , \mathbf{H}_{\Phi^{\ge 2}, (1,1,1,1)} \rangle
  = -q^{2} \neq 0.
\end{equation}

Let $\{\tilde{s}_{\lambda}^{k}\}$ denote the basis dual to $\{s_{\lambda}^k\}$ with respect to the Hall inner product. By \Cref{Teo C}, the expansion of $\tilde{s}_{\lambda}^{k}$ in the basis $\{\tilde{s}_{\mu}^{k+1}\}$ has nonnegative coefficients. It follows, in particular, that $\tilde{s}_{\lambda}^{k}$ is Schur-positive for every $k$ and every $\lambda$.

Blasiak, Morse, and Pun~\cite{blasiak2024demazure} proved the Schur positivity of Catalan symmetric functions for arbitrary root ideals and partitions. Their results provide a powerful refinement of several earlier positivity phenomena. For instance, the Lascoux--Sch\"utzenberger charge formula for $K_{\lambda,\mu}(q)$ appears as a special case of their tableaux formula for Catalan symmetric functions; see \cite[Theorem~2.18]{blasiak2024demazure}.

The functions $\tilde{s}_{\lambda}^{k}$ are more rigid than Catalan symmetric functions, but they have a different advantage: they provide a \emph{global} positive interpolation between modified Hall--Littlewood polynomials and Schur functions. More precisely, an entire basis evolves step by step from the modified Hall--Littlewood basis to the Schur basis. Catalan symmetric functions, by contrast, interpolate one function at a time rather than through a distinguished sequence of full bases.

We also mention the related work~\cite{blasiak2019catalan}, where Blasiak, Morse, Pun, and Summers proved the positivity of the $k$-Schur functions $\mathfrak{s}^k_\lambda$ of Lapointe, Lascoux, and Morse~\cite{luc2003}. They also established the branching property for $k$-Schur functions: the expansion of a $k$-Schur function in the $(k+1)$-Schur basis is positive. Although this statement resembles the positivity satisfied by the dual pre-canonical bases, the two phenomena are different. For $k$-Schur functions, the parameter $k$ controls the allowed size of the first part of the partitions involved; in our setting, the parameter $k$ controls the root ideal used to define the $k$-th pre-canonical Schur function.

We close this discussion by pointing out that readers more familiar with symmetric functions than with Hecke algebras can read much of the paper by applying this dictionary: replace dominant weights by partitions and canonical basis elements by Schur functions.

\subsection{Organization of the paper and acknowledgments}

The paper is organized as follows.
In \Cref{sec: prelimnares}, we introduce the pre-canonical bases and establish the notation used throughout the paper.
In \Cref{section identities to prove first inverse decomposition}, we develop the identities leading to the First Inverse Decomposition, whose proof is given in \Cref{Section: descomposición inversa}.
Likewise, \Cref{sec: Towards the Second Inverse Decomposition} is devoted to developing the tools required for the proof of the Second Inverse Decomposition, which is carried out in \Cref{sec: second inverse decomp}.
Finally, in \Cref{section: positivity}, we prove \Cref{Teo D}, thereby establishing the positivity conjecture for pre-canonical bases. 

\medskip

The authors would like to warmly thank Damian de la Fuente and Nicolas Libedinsky for their comments and suggestions, which greatly improved the presentation of this paper.
The first author would also like to thank Cedric Lecouvey, Cristian Lenart, and Leonardo Patimo for many stimulating discussions related to the topics of this paper during the early stages of the project.

\medskip

David Plaza was partially supported by ANID--FONDECYT Regular Grant 1240199.
Yamil Sagurie was partially supported by ANID-beca doctorado nacional 21201961 and ANID--FONDECYT Regular Grant 1240199.
